% Shared preamble for figures.  Each figure is a standalone .tex in
% figures/Unit_X/ that \input's this file; scripts/make-figures.sh turns it
% into images/Unit_X/<name>.svg.  The lesson page and the guided notes both
% use that one SVG, so a figure is drawn once.
%
% Key figures: a figure whose .tex uses \ifkey ... \fi also gets a
% <name>_KEY.svg, built with \keytrue.  Draw the answers inside \ifkey, in
% the key colour (\ans{} for a label, the "answer" style for lines).
\usepackage[T1]{fontenc}
\usepackage{helvet}
\renewcommand{\familydefault}{\sfdefault}
\usepackage{amsmath}
\usepackage{tikz}
\usetikzlibrary{arrows.meta,calc,angles,quotes,decorations.markings,decorations.pathreplacing}
% axis numbers in the text font, not math italic
\pgfkeys{/pgf/number format/.cd,assume math mode=true,1000 sep={}}

\newif\ifkey
\ifdefined\iskey\keytrue\fi
\definecolor{key}{HTML}{C0192B}      % same red as the guided-notes keys
\newcommand{\ans}[1]{\textcolor{key}{\bfseries #1}}

\tikzset{
  % points: \fill (A) circle[radius=\ptr];
  every picture/.append style={line cap=round, line join=round, >={Stealth[length=2.6mm,width=2mm]}},
  fig/.style={thick},
  line/.style={thick,<->},            % a line goes on forever both ways
  ray/.style={thick,->},              % a ray goes on forever one way
  answer/.style={very thick,key},
  % congruence ticks at the middle of a segment: \draw[tick1] (A) -- (B);
  tick1/.style={postaction={decorate},decoration={markings,
    mark=at position .5 with {\draw[thick] (0,-2.2pt) -- (0,2.2pt);}}},
  tick2/.style={postaction={decorate},decoration={markings,
    mark=at position .5 with {\draw[thick] (-1.3pt,-2.2pt) -- (-1.3pt,2.2pt) (1.3pt,-2.2pt) -- (1.3pt,2.2pt);}}},
  tick3/.style={postaction={decorate},decoration={markings,
    mark=at position .5 with {\draw[thick] (-2.4pt,-2.2pt) -- (-2.4pt,2.2pt) (0,-2.2pt) -- (0,2.2pt) (2.4pt,-2.2pt) -- (2.4pt,2.2pt);}}},
  % parallel arrows at the middle of a line: \draw[par1] (A) -- (B);
  par1/.style={postaction={decorate},decoration={markings,
    mark=at position .5 with {\draw[thick] (-2.5pt,-2.5pt) -- (0,0) -- (-2.5pt,2.5pt);}}},
  par2/.style={postaction={decorate},decoration={markings,
    mark=at position .5 with {\draw[thick] (-4pt,-2.5pt) -- (-1.5pt,0) -- (-4pt,2.5pt) (-1pt,-2.5pt) -- (1.5pt,0) -- (-1pt,2.5pt);}}},
  % angle arcs: \pic[ang] {angle = A--B--C};  ang2 for a second matching angle
  ang/.style={draw,thick,angle radius=5mm},
  ang2/.style={draw,thick,angle radius=5mm,double,double distance=1.2pt},
  % right-angle box: \pic[rt] {right angle = A--B--C};
  rt/.style={draw,thick,angle radius=2.6mm},
}
\newcommand{\ptr}{2pt}
% a labelled point: \pt{A}{(x,y)}{above left}
\newcommand{\pt}[3]{\coordinate (#1) at #2; \fill (#1) circle[radius=\ptr] node[#3] {$#1$};}

% coordinate plane with every unit gridded: \cgrid{xmin}{xmax}{ymin}{ymax}{label step}
% Grid lines every unit. Tick labels every <label step> units, but never closer
% than about 0.55cm: on a small grid a step of 1 becomes 2 (or 5), so the
% numbers do not run into each other. Labels sit on white so grid lines
% do not cross them.
% To leave out tick labels where a curve crosses an axis, before \cgrid:
%   \renewcommand{\cgskipx}{5,-5}  \renewcommand{\cgskipy}{5}
\newcommand{\cgskipx}{}\newcommand{\cgskipy}{}
\makeatletter
\newcommand{\cgrid}[5]{%
  \pgfpointxy{1}{0}\pgfmathsetmacro{\cgunit}{\pgf@x}%
  \pgfmathtruncatemacro{\cgs}{max(#5, (\cgunit*#5 < 15.6) ? 2 : #5, (\cgunit*#5 < 7.9) ? 5 : #5)}%
  \draw[gray!45,thin,step=1] (#1,#3) grid (#2,#4);  % step=1 is one unit; the default is 1cm
  % axes run at least 0.5cm past the last grid line, so the arrowheads clear
  % the end tick labels on a small grid
  \pgfmathsetmacro{\cgo}{max(0.4, 14.2/\cgunit)}%
  \draw[thick,<->] (#1-\cgo,0) -- (#2+\cgo,0) node[right] {$x$};
  \draw[thick,<->] (0,#3-\cgo) -- (0,#4+\cgo) node[above] {$y$};
  \pgfmathtruncatemacro{\cgxa}{ceil((#1)/\cgs)*\cgs}\pgfmathtruncatemacro{\cgxb}{\cgxa+\cgs}%
  \pgfmathtruncatemacro{\cgya}{ceil((#3)/\cgs)*\cgs}\pgfmathtruncatemacro{\cgyb}{\cgya+\cgs}%
  \foreach \t in {\cgxa,\cgxb,...,#2} {\edef\cgt{,\t,}\expandafter\in@\expandafter{\cgt}{,\cgskipx,}\ifin@\else\ifnum\t=0\else\node[below,font=\scriptsize,fill=white,inner sep=1pt,yshift=-1pt] at (\t,0) {\t};\fi\fi}
  \foreach \t in {\cgya,\cgyb,...,#4} {\edef\cgt{,\t,}\expandafter\in@\expandafter{\cgt}{,\cgskipy,}\ifin@\else\ifnum\t=0\else\node[left,font=\scriptsize,fill=white,inner sep=1pt,xshift=-1pt] at (0,\t) {\t};\fi\fi}}
\makeatother
